\scanpage{473}
\mspage{539}

\begin{equation*}
\tag{34}
\left\{
\begin{gathered}
\diagram{\begin{tikzcd}[ampersand replacement=\&]
S\mathcal M_{\rho,\sigma}\arrow[r,"\sim"]\&
\mathcal M_{\rho,\sigma}\mathbin{\underset{\mathbb{G}_m}{\times}}B'_0
\end{tikzcd}}\\
\text{où}\quad\left\{
\begin{gathered}
\diagram{\begin{tikzcd}[ampersand replacement=\&]
\mathcal M_{\rho,\sigma}\arrow[r,"\chi"]\&\mathbb{G}_m\\
B'_0\arrow[r,"{\delta_{B'_0}=\delta_G|B'_0}"]\&\mathbb{G}_m
\end{tikzcd}}
\end{gathered}\right.
\end{gathered}\right.\quad,\quad\text{Ouf !}
\end{equation*}

Cette formule va nous servir de \uline{définition} de
$S\mathcal M_{\rho,\sigma}$ dorénavant -- elle définit tout de suite
une loi de groupe sur $S\mathcal M_{\rho,\sigma}$, et on trouve sur
$S\mathcal M_{\rho,\sigma}$ \uline{deux} structures d'extension

\begin{equation*}
\tag{35}
\diagram{\begin{tikzcd}[ampersand replacement=\&,column sep=small]
1\arrow[r]\&L_0\arrow[r]\arrow[d,phantom,"\Vert"]\&
S\mathcal M_{\rho,\sigma}\arrow[r]\&\mathcal M_{\rho,\sigma}\arrow[r]\&1\,,\\
\&\operatorname{Ker}(B'_0\xrightarrow{\delta_{B'_0}}\mathbb{G}_m)\&
\widetilde x=(x,\unsure{u|y})\arrow[r,mapsto]\&x
\end{tikzcd}}
\end{equation*}

\begin{equation*}
\tag{36}
\diagram{\begin{tikzcd}[ampersand replacement=\&,column sep=small]
1\arrow[r]\&\mathcal M^+_{\rho,\sigma}\arrow[r]\arrow[d,phantom,"\Vert"]\&
S\mathcal M_{\rho,\sigma}\arrow[r]\&B'_0\arrow[r]\&1\\
\&\operatorname{Ker}(\mathcal M_{\rho,\sigma}\xrightarrow{\chi}\mathbb{G}_m)\&
\begin{gathered}\widetilde x\\\Vert\\(x,\underline u)\end{gathered}
\arrow[r,mapsto]\&\underline u
\end{tikzcd}}
\end{equation*}

Le sous-groupe $S\mathcal M_{\rho,\sigma}[0]$ de
$S\mathcal M_{\rho,\sigma}$ correspond à la condition $f=1$ i.e.
$\begin{gathered}\underline U\\\Vert\\\Theta_G(x)\end{gathered}=\underline u$
(qui implique déjà
$\begin{gathered}\det\underline U\\\Vert\\\chi(x)\end{gathered}=\det\underline u$
-- c'est aussi le sous-groupe des $(x,u)$ dans
$\mathcal M_{\rho,\sigma}\times B'_0$ tels que l'on ait
$u=\Theta_G(x)$. Donc $S\mathcal M_{\rho,\sigma}[0]$ est le
sous-groupe de $S\mathcal M_{\rho,\sigma}$ image
\inserted{isomorphique} de $\mathcal M_{\rho,\sigma}[0]$
\deleted{\unsure{par}}
\inserted{(cf (\deleted{\unsure{14|19}}) p. 536')}
par $x\longmapsto(x,\Theta_G(x))$ :

\deleted{\diagram{\begin{tikzcd}[ampersand replacement=\&]
x\arrow[r,mapsto]\&\unsure{(}
\end{tikzcd}}}

\begin{equation*}
\tag{37}
S\mathcal M_{\rho,\sigma}[0]=\operatorname{Im}\left(
\diagram{\begin{tikzcd}[ampersand replacement=\&]
\mathcal M_{\rho,\sigma}[0]\arrow[r,hook]\&S\mathcal M_{\rho,\sigma}\\
x\arrow[r,mapsto]\&(x,\Theta_G(x))
\end{tikzcd}}\right)
\end{equation*}

\[
(\text{où}\quad\mathcal M_{\rho,\sigma}[0]
=\Theta_G^{-1}(B'_0)\quad)
\]

\mspage{472}

\scanpage{474}

La donnée d'un scindage de (35) revient à un scindage de
\[
\diagram{\begin{tikzcd}[ampersand replacement=\&]
1\arrow[r]\&L_0\arrow[r]\&B'_0\arrow[r]\&\mathbb{G}_m\arrow[r]\&1\,,
\end{tikzcd}}
\]
or il y en a un tout prêt par le sous-groupe $T_0$. Donc on en déduit
un plongement de $\mathcal M_{\rho,\sigma}$ dans
$S\mathcal M_{\rho,\sigma}$, qui du coup devient bien un produit
$1/2$-direct.

\begin{equation*}
\tag{38}
S\mathcal M_{\rho,\sigma}\simeq\mathcal M_{\rho,\sigma}\cdot L_0
\end{equation*}

$\mathcal M_{\rho,\sigma}$ \unsure{opérant} d'ailleurs sur $L_0$ via
le caractère $\chi:\mathcal M_{\rho,\sigma}\longrightarrow \mathbb{G}_m$ --
mais cette description me semble maintenant beaucoup moins
``parlante'' que (34), car elle ne met pas en évidence l'homomorphisme
can.
\[
\diagram{\begin{tikzcd}[ampersand replacement=\&]
S\mathcal M_{\rho,\sigma}\arrow[r]\&B'_0
\end{tikzcd}}
\]
qui pourtant est la chose essentielle pour l'introduction de
$S\mathcal M_{\rho,\sigma}$. Une fois ceci posé, on trouve
(``pour mémoire'') le

\uline{Théorème}\quad On a un isom. can.

\begin{equation*}
\tag{39}
\diagram{\begin{tikzcd}[ampersand replacement=\&]
S\mathcal M_{\rho,\sigma}\arrow[r]\&
S\deleted{\unsure{\mathcal M}}\inserted{\unsure{\Sigma}}_{\rho,\sigma}^{*}
\end{tikzcd}}
\quad
\begin{gathered}
\text{schéma en}\\
\text{groupes des}\\
\text{systèmes }(\alpha,\beta,f)\\
\text{satisfaisant }(17)\\
\text{et }(31)
\end{gathered}
\end{equation*}

\deleted{\diagram{\begin{tikzcd}[ampersand replacement=\&]
(x,\underline u)\arrow[r,mapsto]\&
(\alpha=\underline u\unsure{\underline a|\Theta_\rho}^{-1},
\beta=\underline u\underline b^{-1},f=\underline u\underline U^{-1})
\end{tikzcd}}}

\deleted{où $\underline a=\Theta_\rho(x)$,
$\underline b=\Theta_\sigma(x)$, $\underline U=\Theta_G(x)$}

\mspage{473}

\scanpage{475}
\mspage{540}

donné par :

\begin{equation*}
\tag{40}
\begin{gathered}
\diagram{\begin{tikzcd}[ampersand replacement=\&]
(x,\underline u)\arrow[r,mapsto]\&(\alpha,\beta,f)
\end{tikzcd}}
\quad\text{avec}\quad
\alpha=\underline u\underline a^{-1},\quad
\beta=\underline u\underline b^{-1},\quad
f=\underline u\underline U^{-1}\\
\hfill\text{(comme c'est beau !)}\\
\text{où}\quad\underline a=\unsure{\Theta_\rho}(x),\quad
\underline b=\Theta_\sigma(x),\quad\underline U=\Theta_G(x)
\end{gathered}
\end{equation*}

\deleted{\unsure{La} transformation}\inserted{L'opération de $L_0$}
dans $\Sigma_{\rho,\sigma}$ par

\begin{equation*}
\tag{41}
\diagram{\begin{tikzcd}[ampersand replacement=\&]
(\alpha,\beta,f)\arrow[r]\&\unsure{''|(}\ell\alpha,\ell\beta,\ell f
\end{tikzcd}}
\end{equation*}

correspond à l'opération de translation \inserted{à gauche} du
sous-groupe $L_0$ de $S\mathcal M_{\rho,\sigma}$ (cf (35)), formé
des $(1,\ell)$ avec $\ell$ dans $L_0$, donc on récupère
$\mathcal M_{\rho,\sigma}\simeq S\mathcal M_{\rho,\sigma}/L_0$
comme le schéma \inserted{(en groupes)} quotient de
$S\Sigma_{\rho,\sigma}^{*}$ par l'action (41) de $L_0$.

\uline{Remq}\quad Dans cette description des $(\alpha,\beta,f)$
interviennent les quantités nouvelles $\underline u$, $\underline a$,
$\underline b$, $\underline U$, ce qui conduit à considérer le morphisme

\begin{equation*}
\tag{42}
\left\{
\diagram{\begin{tikzcd}[ampersand replacement=\&]
S\mathcal M_{\rho,\sigma}\arrow[r]\&B'_0\times T_\rho\times\unsure{N'_\sigma}\times G\\
(x,\underline u)\arrow[r,mapsto]\&(\underline u,\underline a,\underline b,\underline U)
\end{tikzcd}}\right.
\quad\text{avec}\quad
\begin{gathered}
\underline a=\Theta_\rho(x)\\
\underline b=\Theta_\sigma(x)\\
\underline U=\Theta_G(x)
\end{gathered}
\end{equation*}

On a à ce sujet :

\uline{Proposition}\quad L'hom. (42) est un mono, son image est formée
des systèmes satisfaisant

\begin{equation*}
\tag{43}
\det\underline u=\det\underline b=\det\underline U\quad\text{soit }\mu,
\quad\text{et}\quad\det\underline a/\mu\in\Gamma\mu_2
\end{equation*}

C'est une conséquence immédiate du cor. à la prop de p. 536', et de la
définition de $S\mathcal M_{\rho,\sigma}$ comme produit fibré.

\mspage{474}

\scanpage{476}

10\textsuperscript{o}) Les sous-groupes-sections
$\mathcal M_{\rho,\sigma}(\sigma)$, $\mathcal M_{\rho,\sigma}(\rho)$.

On définit $\mathcal M_{\rho,\sigma}(\sigma)$ par la condition

\[
\beta_0=1\quad\text{i.e.}\quad\beta=f\quad\text{i.e. cf }\unsure{(40)})
\]

\begin{equation*}
\tag{43}
\underline U=\underline b
\end{equation*}

Il est évident par prop p. 536' que le sous-groupe

\begin{equation*}
\tag{44}
\diagram{\begin{tikzcd}[ampersand replacement=\&]
\mathcal M_{\rho,\sigma}(\sigma)\arrow[d,phantom,"\Vert\,\unsure{\mathrm{df}}"]
\subset\mathcal M_{\rho,\sigma})\\
\bigl\{(\underline a,\underline b,\underline U)\in\Gamma\mathcal M_{\rho,\sigma}\mid U=\underline b\bigr\}
\end{tikzcd}}
\end{equation*}

fournit une section -- quand $\underline a,\underline b$ est donné
\inserted{dans $\mathcal T_{\rho,\sigma}$},
\inserted{i.e.} satisfaisant la condition

\[
\unsure{\delta_\rho}(\underline a)/\delta_\sigma(\underline b)\in\Gamma\mu_2,
\]

et la façon la plus simple de le relever en un élément
\inserted{$(\underline a,\underline b,\underline U)$}
de $\mathcal M_{\rho,\sigma}$ n'est-il pas par
$(\underline a,\underline b,\underline U\,\unsure{\mathrm{df}}=\underline b)$ !
(Et on \unsure{aura} bien $\delta_G(\underline U)=\delta_\sigma(\underline b)$ .).
On définit de $\overline{\mathrm m}$ $\mathcal M_{\rho,\sigma}(\rho)$
par $\alpha_0=1$ i.e. $\alpha=f$ i.e.

\begin{equation*}
\tag{45}
\diagram{\begin{tikzcd}[ampersand replacement=\&]
\mathcal M_{\rho,\sigma}(\rho)\arrow[d,phantom,"\Vert"]\subset\mathcal M_{\rho,\sigma}\\
\bigl\{(\underline a,\underline b,\underline U)\in\mathcal M_{\rho,\sigma}\mid\underline U=\underline a\bigr\}
\end{tikzcd}}
\end{equation*}

ce qui fait une section partielle au \inserted{dessus} du sous-groupe
$\mathcal T'_{\rho,\sigma}$ de $\mathcal T_{\rho,\sigma}$ défini comme

\begin{equation*}
\tag{46}
\diagram{\begin{tikzcd}[ampersand replacement=\&]
\mathcal T'_{\rho,\sigma}\arrow[d,phantom,"\Vert\,\unsure{\mathrm{df}}"]\unsure{\in|\subset}\mathcal T_{\rho,\sigma}\\
\bigl((\underline a,\underline b)\in T_\rho\times N'_\sigma\mid\unsure{\delta_\rho}(\underline a)=\delta_\sigma(\underline b)\bigr)
\arrow[d,phantom,"\Vert"]\\
\operatorname{Ker}\bigl(\mathcal T_{\rho,\sigma}\xrightarrow{\epsilon_3}\mu_2\bigr)
\end{tikzcd}}
\end{equation*}

donc

\begin{equation*}
\tag{47}
\mathcal T'_{\rho,\sigma}=T_\rho\times_{\mathbb{G}_m}\unsure{N'_\sigma},\quad\text{et on a}
\end{equation*}

\mspage{475}

\scanpage{477}
\mspage{541}

donc

\begin{equation*}
\tag{48}
\diagram{\begin{tikzcd}[ampersand replacement=\&]
\mathcal T'_{\rho,\sigma}/\mathcal T^0_{\rho,\sigma}\arrow[r,"\epsilon_2"]\&\mu_2
\end{tikzcd}}\,.
\end{equation*}

\uline{Remq}\quad Ces sous-groupes peuvent se définir, et fournir une
section resp une section \deleted{sur}\inserted{au-dessus de}
$\mathcal T'_{\rho,\sigma}$, dès qu'on a (dans la situation générale
avec $\unsure{\mathfrak f|\mathfrak G}$, p. 526' ff) des plongements
\diagram{\begin{tikzcd}[ampersand replacement=\&]
T_\rho\arrow[r,hook]\&G
\end{tikzcd}},
\diagram{\begin{tikzcd}[ampersand replacement=\&]
N'_\sigma\arrow[r,hook]\&G
\end{tikzcd}}, tels que
$\delta_\rho,\delta_\sigma$ soient induits par $\delta_G$.

Dans le cas modulaire, on écrit aussi

\begin{equation*}
\tag{49}
\mathcal M_{\rho,\sigma}(1/2)=\mathcal M_{\rho,\sigma}(\sigma),\quad
\mathcal M_{\rho,\sigma}(-\bar j)=\mathcal M_{\rho,\sigma}(\rho),
\end{equation*}

et on définit plus \unsure{général\textsuperscript{t}|généralement}

\begin{equation*}
\tag{50}
\mathcal M_{\rho,\sigma}(q),\qquad
q\in\{\unsure{0},1,\infty\};1/2,\unsure{\underline 2|2},-1\}
\end{equation*}

à partir de $\mathcal M_{\rho,\sigma}(0)$ et
$\mathcal M_{\rho,\sigma}(1/2)$ par

\begin{equation*}
\tag{51}
\operatorname{int}(\rho_G)\mathcal M_{\rho,\sigma}(q)
=\mathcal M_{\rho,\sigma}(\dot\rho(q))
\end{equation*}

\margin{$\rho_G$ est $\rho$ en tant qu'élément du sous-groupe
$\unsure{\mathcal G}_{\rho,\sigma}^{+}\simeq SG$ de
$\underset{\unsure{\sim}}{\mathcal M}_{\rho,\sigma}$}

$\dot\rho\deleted{(q)}$ est la transformation homographique

\begin{equation*}
\tag{52}
\dot\rho:\quad
\diagram{\begin{tikzcd}[ampersand replacement=\&]
q\arrow[r,mapsto]\&\dfrac{1}{1-q}
\end{tikzcd}}
\end{equation*}

(dont l'itéré troisième est l'identité, tout comme l'itéré troisième
de $\widetilde\rho$).

Je \unsure{voudrais|voudrait} : définir $\mathcal M_{\rho,\sigma}(-j)$, et
\deleted{\unsure{comme}} on a le choix à l'intérieur d'une classe de
$\underline{\unsure{\mathrm{Sl}|\mathrm{SL}}(2,\mathbb Z)}^{\wedge}$-conjugaison bien déterminée,
un choix plus

\mspage{476.}

\scanpage{478}
\unsure{restreint} serait entre
\[
\operatorname{int}(\unsure{\sigma_i})(\mathcal M_{\rho,\sigma}(0))
\qquad i\in\{0,1,\infty\}
\tag{53}
\]
\mspage{477}

\scanpage{479}
\mspage{542}
\uline{$\overline{\mathrm{XII}}$}\qquad\uline{Rétrospective}

Par bien des détours inutilement compliqués, j'ai fini de façon bien simple par construire deux schémas en groupes sur $\mathbb Z$
\[
\mathcal M_{\rho,\sigma}\subset
\deleted{\underset{\substack{\Vert\\\unsure{\mathrm{Gl}(2)}}}{G}\times}
T_\rho\times N'_\sigma\times G
\tag{1}
\]
\[
S\mathcal M_{\rho,\sigma}\subset
\unsure{B'_0}\times T_\rho\times N'_\sigma\times
\unsure{\deleted{G}B_0}
\tag{2}
\]
par
\[
\mathcal M_{\rho,\sigma}=
\left\{\overbrace{\left(
\unsure{\deleted{\underline a}\underline U},
\unsure{\deleted{\underline b}\underline a},
\unsure{\deleted{\underline U}\underline b}
\right)}^{x}\,\middle|\,
\overbrace{\det\underline U=\det\underline b}^{\chi(x)}=
\overset{\substack{\varepsilon_3(x)\\\vert}}{\varepsilon'}\det\underline a,
{\varepsilon'}^2=1\right\}
\tag{3}
\]
\[
S\mathcal M_{\rho,\sigma}=
\left\{(\unsure{\underline u},\underline a,\underline b,\underline U)\,\middle|\,
\det\underline U=\det\underline u=\det\underline b=
\varepsilon'\det\underline a,\ {\varepsilon'}^2=1\right\}
\tag{4}
\]
\margin{donc $S\mathcal M_{\rho,\sigma}=B'_0\times_{\mathbb{G}_m}\mathcal M_{\rho,\sigma}$}

avec
\[
\mathcal M_{\rho,\sigma}\simeq S\mathcal M_{\rho,\sigma}/L_0
\tag{5}
\]
où
\diagram{\begin{tikzcd}[ampersand replacement=\&,column sep=large]
L_0\arrow[r,hook]\&S\mathcal M_{\rho,\sigma}\\
\lambda\arrow[r,mapsto]\&(1,1,1,\lambda)
\end{tikzcd}}

et
\[
S\mathcal M_{\rho,\sigma}\simeq
\underset{\text{produit }\frac12\text{ direct}}{\mathcal M_{\rho,\sigma}\mathbin{.}L_0}
\tag{6}
\]
où
\[
\tag{7}
\left\{\vcenter{\hbox{\diagram{\begin{tikzcd}[ampersand replacement=\&,column sep=large]
\mathcal M_{\rho,\sigma}\arrow[r,hook]\&S\mathcal M_{\rho,\sigma}\\
x=(\underline U,\underline a,\underline b)\arrow[r,mapsto]\&
(\underline U,\underline a,\underline b,
 i_{T_0}(\unsure{\deleted{\det}\mu})\deleted{\unsure{(\underline U)}})
\end{tikzcd}}}}\right.
\qquad\text{où }\mu=\chi(x).
\]

Ici (pour mémoire) ds le cas modulaire
\[
\rho=\begin{pmatrix}0&1\\-1&1\end{pmatrix}
\qquad
\sigma=\begin{pmatrix}0&1\\-1&0\end{pmatrix}
\tag{8}
\]
\[
\left\{\begin{aligned}
T_\rho&=\left\{(c\rho+d)\,\middle|\,
\underbrace{\det(c\rho+d)}_{c^2+cd+d^2}\ \mathrm{inv.}\right\}\\
T'_\sigma&=\left\{(t\sigma+u)\,\middle|\,
\underbrace{\det(t\sigma+u)}_{t^2+u^2}\ \mathrm{inv.}\right\}.
\end{aligned}\right.
\tag{9}
\]
\margin{$\left\{\begin{aligned}
B'_0&=\left\{\begin{pmatrix}\mu&\lambda\\0&1\end{pmatrix}\,\middle|\,\mu\ \mathrm{inv.}\right\}\\
T_0&=\left\{\begin{pmatrix}\mu&0\\0&1\end{pmatrix}\,\middle|\,\mu\ \mathrm{inv.}\right\}\\
i_{T_0}(\mu)&=\begin{pmatrix}\mu&0\\0&1\end{pmatrix}\\
L_0&=\left\{\begin{pmatrix}1&\lambda\\0&1\end{pmatrix}\right\}
\end{aligned}\right.$}
\deleted{(\unsure{10})}
\[
N'_\sigma=\left\{\underline b\text{ dans }G\,\middle|\,
\underline b(\sigma)=\underset{\substack{\Vert\\\unsure{\varepsilon_\sigma}(\underline b)}}{\varepsilon}\sigma,
\ \varepsilon^2=1\right\},
\tag{10}
\]
d'où suite exacte
\[
\tag{11}
\vcenter{\hbox{\diagram{\begin{tikzcd}[ampersand replacement=\&,column sep=large]
1\arrow[r,dash]\&T_\sigma\arrow[r]\&N'_\sigma\arrow[r,"{\unsure{\varepsilon_\sigma}}"]\&\mu_2\arrow[r]\&1
\end{tikzcd}}}},
\]
\mspage{478}

\scanpage{480}
et les trois hom.
\[
\tag{12}
\vcenter{\hbox{\diagram{\begin{tikzcd}[ampersand replacement=\&,column sep=large]
\mathcal M_{\rho,\sigma}\arrow[r,"\Theta_G"]\&G
\end{tikzcd}}}},\quad
\vcenter{\hbox{\diagram{\begin{tikzcd}[ampersand replacement=\&,column sep=large]
\mathcal M_{\rho,\sigma}\arrow[r,"\Theta_\rho"]\&T_\rho
\end{tikzcd}}}},\quad
\vcenter{\hbox{\diagram{\begin{tikzcd}[ampersand replacement=\&,column sep=large]
\mathcal M_{\rho,\sigma}\arrow[r,"\Theta_\sigma"]\&N'_\sigma
\end{tikzcd}}}}
\]
donnés dans la définition (1)(3) se complètent par
\[
\tag{12 bis}
\vcenter{\hbox{\diagram{\begin{tikzcd}[ampersand replacement=\&,column sep=large]
\mathcal M_{\rho,\sigma}\arrow[r,"{\chi=\det\circ\Theta_G}"]\&\mathbb{G}_m
\end{tikzcd}}}},\quad
\vcenter{\hbox{\diagram{\begin{tikzcd}[ampersand replacement=\&,column sep=large]
\mathcal M_{\rho,\sigma}\arrow[r,"{\varepsilon_2=\varepsilon_\sigma\circ\Theta_\sigma}"]\&\mu_2
\end{tikzcd}}}},\quad
\vcenter{\hbox{\diagram{\begin{tikzcd}[ampersand replacement=\&,column sep=large]
\mathcal M_{\rho,\sigma}\arrow[r,"{\varepsilon_3(x)=\frac{\det\underline a}{\det\underline b}}"]\&\mu_2
\end{tikzcd}}}}
\]
de sorte qu'on trouve
\[
\tag{13}
\vcenter{\hbox{\diagram{\begin{tikzcd}[ampersand replacement=\&,column sep=large]
\mathcal M_{\rho,\sigma}\arrow[r,hook,"{(\Theta_G,\Theta_\rho,\Theta_\sigma,\chi,\varepsilon_2,\varepsilon_3)}"]\&
G\times T_\rho\times N'_\sigma\times\mathbb{G}_m\times\mu_2\times\mu_2\\
x\arrow[r,mapsto]\&(\underline U,\underline a,\underline b,\mu,\varepsilon,\varepsilon')
\end{tikzcd}}}}
\]
dont l'image est formée des systèmes \unsure{satisfaisant}
\[
\left\{\begin{aligned}
&\det\underline U=\det\underline b=\varepsilon'\det\underline a=\mu\\
&\varepsilon=\varepsilon_\sigma(\underline b)\quad\text{i.e}\quad
\begin{gathered}
\underline b(\sigma)=\varepsilon\sigma\quad\text{i.e.}\\
\underbrace{\sigma(\underline b)}_{{}^tb^{-1}}=\varepsilon\underline b
\quad\text{i.e}\quad \underline b\,{}^tb=\varepsilon
\end{gathered}
\end{aligned}\right.
\tag{14}
\]

De même on a mono
\[
\tag{15}
\vcenter{\hbox{\diagram{\begin{tikzcd}[ampersand replacement=\&,column sep=large]
S\mathcal M_{\rho,\sigma}\arrow[r,hook]\&
\unsure{B'_0}\times T_\sigma\times N'_\sigma\times
\underline{\unsure{\deleted{G}B_0}}\times\mathbb{G}_m\times\mu_2\times\mu_2\\
x\arrow[r,mapsto]\&(\unsure{\underline u},\underline a,\underline b,\underline U,\mu,\varepsilon,\varepsilon')
\end{tikzcd}}}}
\]
dont l'image est formée des systèmes \unsure{satisfaisant}
\[
\left\{\begin{aligned}
\mu&=\det\underline u=\det\underline U=\det\underline b=\varepsilon'\det\underline a\\
\varepsilon&=\varepsilon_\sigma(\underline b)\quad\text{i.e}\qquad
\underline b(\sigma)=\varepsilon(\sigma)
\end{aligned}\right.
\tag{16}
\]

\inserted{Soit} d'autre part
\[
\tag{17}
\begin{gathered}
\vcenter{\hbox{\diagram{\begin{tikzcd}[ampersand replacement=\&,column sep=large]
S\Sigma^*_{\rho,\sigma}\arrow[r,hook]\&
G\times\deleted{\unsure{G}}G\times G\times\mathbb G_a\times\mu_2
\end{tikzcd}}}}\\[-2pt]
\Vert\\[-2pt]
\left\{\underbrace{(\alpha,\beta,f,\mu,\varepsilon)}_y\,\middle|\,
\begin{gathered}
(\det\alpha)^2=\det\beta=\det f=1\\
\underbrace{\alpha\rho(\alpha)^{-1}}_\xi=
\underbrace{\varepsilon\beta\sigma(\beta)^{-1}}_{\unsure{\eta}}
\varepsilon_1(\mu-1)
\end{gathered}\right\}
\end{gathered}
\]
(où on définit
\[
\varepsilon_1=\begin{pmatrix}1&0\\1&1\end{pmatrix},\qquad
\varepsilon_1(\nu)=\varepsilon_1^\nu\overset{\unsure{\mathrm{dfn}}}{=}
\begin{pmatrix}1&0\\\nu&1\end{pmatrix}.)
\tag{18}
\]
\mspage{479}

\scanpage{481}
\mspage{543}
en fait la projection
\[
\tag{19}
\left\{\vcenter{\hbox{\diagram{\begin{tikzcd}[ampersand replacement=\&,column sep=large]
S\Sigma^*_{\rho,\sigma}\arrow[r,hook]\&G\times G\times G\\
y\arrow[r,mapsto]\&(\alpha,\beta,f)
\end{tikzcd}}}}\right.
\]
est déjà un mono. On fait opérer \deleted{$L_0$}\inserted{librement}
\[
L_0=\operatorname{Ker}\left(
\vcenter{\hbox{\diagram{\begin{tikzcd}[ampersand replacement=\&,column sep=large]
B'_0\arrow[r,"\det"]\&\mathbb{G}_m
\end{tikzcd}}}}\right)
=\deleted{\unsure{\{}}\left\{\begin{pmatrix}1&\lambda\\0&1\end{pmatrix}\right\}
\quad\text{sur }S\Sigma^*_{\rho,\sigma}
\]
par
\[
\tag{20}
\vcenter{\hbox{\diagram{\begin{tikzcd}[ampersand replacement=\&,column sep=large]
(\alpha,\beta,f)\arrow[r,mapsto]\&(l\alpha,l\beta,lf)
\end{tikzcd}}}},
\]
et on pose
\[
\Sigma^*_{\rho,\sigma}\overset{\unsure{\mathrm{dfn}}}{=}
\deleted{S\Sigma^*_{\rho,\sigma}}\ L_0\backslash S\Sigma^*_{\rho,\sigma}.
\tag{21}
\]

Ceci \unsure{posé, on a} un \uline{isomorphisme}
\[
\tag{22}
\vcenter{\hbox{\diagram{\begin{tikzcd}[ampersand replacement=\&,column sep=large]
S\mathcal M_{\rho,\sigma}\arrow[r,"\sim"]\&S\Sigma^*_{\rho,\sigma}\\
(\underline u,\underline a,\underline b,\underline U)\arrow[r,mapsto]\&
\left\{\begin{gathered}
(\alpha,\beta,f),\\
\alpha=\underline u\,\underline a^{-1},\quad
\beta=\underline u\,\underline b^{-1},\quad
f=\underline u\,\underline U^{-1}
\end{gathered}\right.
\end{tikzcd}}}}
\]
transformant l'action de $L_0$ sur $S\unsure{\deleted{\Sigma^*}\mathcal M}_{\rho,\sigma}$ \inserted{par \unsure{transl.} à g.} \inserted{en} l'action de $L_0$ décrite dans (20), et induisant un iso
\[
\tag{23}
\vcenter{\hbox{\diagram{\begin{tikzcd}[ampersand replacement=\&,column sep=large]
\mathcal M_{\rho,\sigma}\arrow[r,"\sim"]\&\Sigma^*_{\rho,\sigma}
\end{tikzcd}}}}
\]
par passage au quotient.

On a des hom. canoniques
\[
\tag{24}
\vcenter{\hbox{\diagram{\begin{tikzcd}[ampersand replacement=\&,column sep=large]
S_0\mathcal M_{0,3}^{\sim}\arrow[r]\&S\mathcal M_{\rho,\sigma}(\widehat{\mathbb Z})/\Sigma
\end{tikzcd}}}}
\]
\[
\tag{25}
\vcenter{\hbox{\diagram{\begin{tikzcd}[ampersand replacement=\&,column sep=large]
\mathcal M_{0,3}^{\sim}\arrow[r]\&\mathcal M_{\rho,\sigma}(\widehat{\mathbb Z})/\Sigma
\end{tikzcd}}}}
\]
le deuxième déduit du premier par passage au quotient, et coïncidant sur $\mathcal G_{0,3}^{+\wedge}$ avec
\deleted{\unsure{et|le} \illegible{2} \unsure{en groupes}}
\mspage{480}

\scanpage{482}
l'hom composé
\[
\tag{26}
\vcenter{\hbox{\diagram{\begin{tikzcd}[ampersand replacement=\&,column sep=large]
\mathcal G_{0,3}^{+\wedge}
\arrow[r,dash,"{\substack{\S42\ \overline{\underline{\mathrm{IV}}}\\\sim}}"]\&
\underset{\substack{\Vert\\SG(\widehat{\mathbb Z})/SG(\widehat{\mathbb Z})\cap\Sigma}}
{\mathrm{Sl}(2,\widehat{\mathbb Z})'}
\arrow[r,hook]\&
\deleted{\unsure{\mathrm{Sl}(2,\widehat{\mathbb Z})'}}
\mathcal M_{\rho,\sigma}(\widehat{\mathbb Z})/\Sigma
\end{tikzcd}}}}
\]
grâce à l'hom canonique
\[
\tag{27}
\vcenter{\hbox{\diagram{\begin{tikzcd}[ampersand replacement=\&,column sep=large,row sep=small]
SG\arrow[r,hook]\arrow[d,equal]\&\mathcal M_{\rho,\sigma}\\
\mathrm{Sl}(2)\deleted{\unsure{\mathbb Z}}\arrow[r,equal]\&
\left\{(\unsure{\underline a},\underline b,\underline U)
\,\middle|\,\overset{\inserted{\text{dans }\mathcal M_{\rho,\sigma}}}
{\underline a=\underline b=1}\right\}
\end{tikzcd}}}}
\]

On a défini\deleted{\unsure{t}}
\[
\tag{28}
\vcenter{\hbox{\diagram{\begin{tikzcd}[ampersand replacement=\&,column sep=large]
\Sigma=\underset{\substack{\Vert\\(\mu_2)^3(\mathbb Z)}}{(\mathbb Z/2\mathbb Z)^3}
\arrow[r,hook]\&\mathcal M_{\rho,\sigma}(\mathbb Z)
\arrow[r,phantom,"\subset"]\&\mathcal M_{\rho,\sigma}(\widehat{\mathbb Z})
\end{tikzcd}}}}
\]
via l'hom canonique (central)
\[
\tag{29}
\vcenter{\hbox{\diagram{\begin{tikzcd}[ampersand replacement=\&,column sep=large]
\mu_2^3\arrow[r]\&\deleted{\unsure{\mathcal T}}\mathcal M_{\rho,\sigma}
\arrow[r,phantom,"\subset"]\&T_\rho\times N'_\sigma\times G
\end{tikzcd}}}}
\]
défini via les composés canoniques
\[
\tag{30}
\vcenter{\hbox{\diagram{\begin{tikzcd}[ampersand replacement=\&,column sep=large,row sep=small]
\&\&T_\rho\&\\
\mu_2\arrow[r,hook]\&\mathbb{G}_m\arrow[ur,hook]\arrow[r,hook]\arrow[dr,hook]\&
T_\sigma\arrow[r,hook]\&N'_\sigma\\
\&\&G\&
\end{tikzcd}}}}.
\]

Par passage au quotient, l'hom (25) induit
\[
\tag{\unsure{31|34}}
\vcenter{\hbox{\diagram{\begin{tikzcd}[ampersand replacement=\&,column sep=large]
\Gamma_{\deleted{\unsure{0}}0,3}^{\sim}\arrow[r]\&
\mathcal T_{\rho,\sigma}(\widehat{\mathbb Z})/
\underbrace{\{\pm1\}\times\{\pm1\}}_{\Sigma_0}
\end{tikzcd}}}},
\]
où
\[
\tag{32}
\begin{gathered}
\vcenter{\hbox{\diagram{\begin{tikzcd}[ampersand replacement=\&,column sep=large]
\mathcal T_{\rho,\sigma}\arrow[r,hook]\&T_\rho\times N'_\sigma
\end{tikzcd}}}}\\
\Vert\\
\left\{(\underline a,\underline b)\,\middle|\,
\det\underline a/\det\underline b\in\Gamma\mu_2\right\}
\end{gathered}
\]

La donnée --- plutôt la \uline{construction} faite par nous --- d'un \deleted{(24)} hom de groupes (24)
\mspage{481}

\scanpage{483}
\mspage{544}

revient à la construction des quatre homomorphismes de groupes
\[
\diagram{\begin{tikzcd}[ampersand replacement=\&,column sep=huge]
S_0\mathcal M_{0,3}^{\sim}
  \arrow[r, "{\substack{\text{donnant }\underline u\\\Theta_{B'_0}^{\pi}}}"]
  \& B'_0(\widehat{\mathbb Z})=\widehat{\mathbb Z}^{*}\mathbin{.}\widehat{\mathbb Z}
\end{tikzcd}}
\qquad\text{(produit }\tfrac12\text{ direct)}
\tag{33}
\]
\[
\diagram{\begin{tikzcd}[ampersand replacement=\&,column sep=huge]
\mathcal M_{0,3}^{\sim}
  \arrow[r, "{\substack{\text{donnant }\underline U\text{ (mod signe)}\\\Theta_G^{\pi}}}"]
  \&G(\widehat{\mathbb Z})/\pm1=\mathrm{Gl}(2,\widehat{\mathbb Z})/\pm1
\end{tikzcd}}
\tag{34}
\]
\[
\diagram{\begin{tikzcd}[ampersand replacement=\&,column sep=huge]
\Gamma_{0,3}^{\sim}
  \arrow[r, "{\substack{\text{donnant }\underline a\text{ (mod signe)}\\\Theta_{\rho}^{\pi}}}"]
  \&T_{\rho}(\widehat{\mathbb Z})/\pm1
\end{tikzcd}}
\quad\begin{gathered}
\text{(unités adéliques}\\
\text{(mod }\pm1\text{) du}\\
\text{corps }\mathbb Q(\deleted{\unsure{\sqrt[3]{1}}}\inserted{j})\text{)}
\end{gathered}
\tag{35}
\]
\[
\diagram{\begin{tikzcd}[ampersand replacement=\&,column sep=huge]
\Gamma_{0,3}^{\sim}
  \arrow[r, "{\substack{\text{donnant }\underline b\text{ (mod signe)}\\\Theta_{\sigma}^{\pi}}}"]
  \&T_{\sigma}(\widehat{\mathbb Z})/\pm1
\end{tikzcd}}
\quad\begin{gathered}
\text{(unités adéliques}\\
\text{(mod }\pm1\text{) du}\\
\text{corps }\mathbb Q(i)\text{)}
\end{gathered}
\tag{36}
\]
\unsure{reliés|déliés} par les relations (4), que nous allons identifier.
On utilise le sous-groupe
\[
\begin{aligned}
\mathcal M_{0,3}^{!\sim}[0]^{\unsure{-}}
  &=\text{sous-groupe de }\mathcal M_{0,3}^{!\sim}\text{ qui}\\
  &\quad\text{(invarie les classes des \unsure{s-groupes}}\\
  &\quad\text{lacets de }\Pi_{0,3}^{\wedge}\text{, et qui)}\\
  &\quad\text{normalise }L_{0\pi}\text{ (sous-groupe}\\
  &\quad\text{fermé engendré par }l_0\text{)}
\end{aligned}
\tag{37}
\]
et on note que
\[
\diagram{\begin{tikzcd}[ampersand replacement=\&]
1\arrow[r]\&L_{0\pi}\arrow[r]\&\mathcal M_{0,3}^{!\sim}[0]\arrow[r]\&\Gamma_{0,3}^{\sim}\arrow[r]\&1
\end{tikzcd}}
\tag{38}
\]
et que
\[
\left\{\begin{aligned}
\mathcal M_{0,3}^{\sim}
  &=\mathcal M_{0,3}^{!\sim}[0]\mathbin{.}
  \underbrace{\mathcal G_{0,3}^{+\wedge}}_{\substack{\text{sous-groupe invariant}\\\inserted{\text{de }\mathcal M_{0,3}^{\sim}}}}\\
\text{avec}\qquad
\mathcal M_{0,3}^{!\sim}[0]\cap\mathcal G_{0,3}^{+\wedge}
  &=L_{0\pi}
\end{aligned}\right.
\tag{39}
\]
--- bien sûr (38) est conséquence de (39).

enfin
\[
S_0\mathcal M_{0,3}^{\sim}\simeq
\underbrace{\mathcal M_{0,3}^{!\sim}[0]\mathbin{.}\mathcal G_{0,3}^{+\wedge}}_{\text{(produit }\tfrac12\text{ direct)}}
\tag{40}
\]
1\textsuperscript{o}) L'hom. (33) \inserted{de $S_0\mathcal M_{0,3}^{\sim}$}
\deleted{\unsure{c'}}est trivial sur le facteur $\mathcal G_{0,3}^{+\wedge}$

\mspage{482}

\scanpage{484}

\margin{donc la \deleted{définition}\inserted{description}
de (33) revient à la description du \unsure{seul hom.|premier hom.}
\diagram{\begin{tikzcd}[ampersand replacement=\&]
\mathcal M_{0,3}^{\sim}[0]\arrow[r]\&B'_0(\widehat{\mathbb Z})
\end{tikzcd}}}
Sa \deleted{\unsure{ne}} \unsure{\uline{valeur}} ne dépend que du couple
$(\alpha,\beta)$ \deleted{\unsure{part|pris}} du triplet $(\alpha,\beta,f)$
qui décrit un élément de $\mathcal M_{0,3}^{\sim}$), et
la \deleted{\unsure{posi}} définition de $\underline u$ en termes de
$\alpha,\beta$ --- et $\overline{\mathrm m}$ seulement en termes
de l'image de $\xi=\alpha\rho(\alpha)^{-1}$ dans
$\mathrm{Sl}(2,\widehat{\mathbb Z})$
\deleted{\unsure{mod} $\unsure{\pm1}$, \unsure{ou}}
\inserted{et $\overline{\mathrm m}$} seulement dans
$\mathrm{Sl}(2,\widehat{\mathbb Z})'$ ! ---
\margin{Car on \unsure{normalise} par la condition
$\xi\rho(\xi)\rho^2(\xi)=\unsure{\pm1|+1}$ !}
\unsure{à|a} été \unsure{fait|faite} essentiellement dans le \S45
ou tout au moins dans \S46 $\overline{\underline{\mathrm{III}}}$, --- on a
\[
\underline u:=\begin{pmatrix}1/D&B/D\\0&1\end{pmatrix}
\qquad\begin{gathered}
\text{où }\xi=\begin{pmatrix}A&B\\C&D\end{pmatrix}\\
\text{(dans }\mathrm{Sl}(2,\widehat{\mathbb Z})\text{)}
\end{gathered}
\tag{41}
\]
ou encore, en termes de
\[
\underline\alpha=\begin{pmatrix}a&b\\c&d\end{pmatrix}
\qquad\text{avec }ad-bc\unsure{=|\in}\inserted{\varepsilon'\in}\{\pm1\}
\]
\[
\left\{\begin{aligned}
\inserted{\varepsilon'}D&=(c^2+cd+d^2)\deleted{\unsure{/\varepsilon}}\\
\inserted{\varepsilon'}B&=ac+bd+bc
\end{aligned}\right.
\qquad\deleted{\text{avec }\unsure{\varepsilon'}}
\tag{42}
\]

Cela ne dépend pas des fondements non faits sur les groupes profinis ---
les ``calculs en $\xi,\eta$'' sont en fait suffisants.
C'est grâce à la \unsure{description} de cet homomorphisme
qu'on peut définir
\[
\begin{gathered}
\mathcal M_{0,3}^{\sim}(0)\subset\mathcal M_{0,3}^{\sim}[0]\\
\begin{array}{c}
\Vert\\
\{\underline{\underline u}\mid\underline u\in T_0(\widehat{\mathbb Z})
\quad\text{i.e. }B=0\}\\
\qquad\text{i.e. }ac+bd+bc=0
\end{array}
\end{gathered}
\tag{43}
\]
(condition de normalisation),

\mspage{483}

\scanpage{485}
\mspage{545}

et obtenir des décompositions canoniques en produits
$\tfrac12$-\unsure{directs|direct}
\[
\mathcal M_{0,3}^{\sim}[0]\simeq\mathcal M_{0,3}^{\sim}(0)\mathbin{.}L_{0\pi}
\tag{44}
\]
\[
\mathcal M_{0,3}^{\sim}\simeq\mathcal M_{0,3}^{\sim}(0)\mathbin{.}\mathcal G_{0,3}^{+\wedge}.
\tag{45}
\]

2\textsuperscript{o}) Pour décrire (34), comme il est déjà connu sur
$\mathcal G_{0,3}^{+\wedge}$, il suffit de le décrire sur
$\mathcal M_{0,3}^{\sim}[0]$ ou $\overline{\mathrm m}$ seulement sur
$\mathcal M_{0,3}^{\sim}(0)$ (moyennant bien sûr des compatibilités\ldots).
Et sur $\mathcal M_{0,3}^{\sim}[0]$ il coïncide avec (33) ---
tandis que sur $\mathcal M_{0,3}^{\sim}(0)$ il est pratiquement trivial
et se réduit au multiplicateur $\mu$
\[
\Theta_G^{\pi}=\Theta_{B'_0}^{\pi}
\quad\text{sur }\mathcal M_{\unsure{\rho|0},3}^{\sim}[0]
\tag{46}
\]
\[
\left\{\begin{aligned}
\Theta_G^{\pi}(\underline{\underline u})
  &=i'_{T_0}(\mu)
  \qquad\begin{gathered}\text{où }\mu\text{ est le}\\
  \text{multiplicateur de }\underline{\underline u}\end{gathered}\\
\text{\unsure{si|où}}\quad
\underline{\underline u}&\in\mathcal M_{0,3}^{\sim}(0)
\end{aligned}\right.
\tag{47}
\]
(c'est $\overline{\mathrm m}$ une définition possible de $\mathcal M_{0,3}^{\sim}(0)$).

\uline{Remarque} Comme $\Theta_G^{\pi}$ a été construit
``à la force des \unsure{poignet|poignets}'' par l'intermédiaire
de $\Theta_{B'_0}^{\pi}$, i.e. $\underline U$ a été construit par $\underline u$,
il y a eu longtemps une espèce de

\mspage{484}

\scanpage{486}

confusion dans mon esprit entre $\underline u$ et $\underline U$,
d'autant plus que la plus grande partie du temps je me réduisais à travailler
dans $\mathcal M_{\unsure{0|\rho},3}^{\sim}[0]$
(après avoir pas mal cafouillé avec des points de vue plus
\unsure{variés}, avant de me rendre compte que techniquement c'était
la façon la plus élégante de procéder --- à partir du \S\unsure{46|45}).
Cette confusion ne s'est définitivement clarifiée que hier et aujourd'hui,
à force de mettre et remettre les pieds dans mes propres plats\ldots

3\textsuperscript{o}) Les hom. (35)(36) n'ont été obtenus
qu'en travaillant à fond sur les paramètres $\alpha,\beta$
\unsure{pour} $\underline u$ --- ici aussi, il suffit de travailler
dans $\mathcal M_{0,3}^{\sim}[0]$, et de définir
\[
\diagram{\begin{tikzcd}[ampersand replacement=\&,row sep=tiny,column sep=huge]
\underline{\underline u}\arrow[r,mapsto]\&\underline a\\
\mathcal M_{0,3}^{\sim}[0]\arrow[r]\&T_{\rho}(\widehat{\mathbb Z})/\pm1\\
{}\arrow[r]\&\unsure{N'_{\sigma}}(\widehat{\mathbb Z})/\pm1\\
\underline{\underline u}\arrow[r,mapsto]\&\underline b
\end{tikzcd}}
\]

\mspage{485}

\scanpage{487}
\mspage{546}

de façon que ce soit trivial sur $L_0^{\deleted{\unsure{\xi|\wedge}}\pi}$.

Les quantités $\underline a,\underline b$ sont apparues dans les
calculs progressivement à partir du
\S44 (p. 366$'$ ff), la lumière sur leur
rôle a commencé à se faire seulement
à partir du \S48 $\overline{\underline{\mathrm{III}}},
\overline{\underline{\mathrm{IV}}},\overline{\underline{\mathrm{V}}}$ ---
et une compréhension complète de leur nature, du point
de vue des schémas en groupes qui interviennent dans la théorie,
s'est faite également seulement en ces \unsure{tous|tout} derniers
jours. Cette compréhension reste d'ailleurs pour le moment assez
formelle, et il faudrait approfondir
dans plusieurs directions,

a) Expliciter la signification du point
de vue des corps de classes de
$\mathbb Q(j)$ et \inserted{de} $\mathbb Q(\deleted{j}\inserted{i})$.

b) Expliciter la signification pour
la théorie des cartes ``\unsure{congruentielles|congruencielles}''
(i.e. \unsure{subordonnées|subordonnés} à des repr. \deleted{\unsure{contin.}}
du quotient $\mathrm{Sl}(2,\widehat{\mathbb Z})'$ de
$\mathrm{Sl}(2,\mathbb Z)^{\wedge\prime}$)
des résultats obtenus,

\mspage{486}

\scanpage{488}

c) Introduire sur $\mathcal M_{\rho,\sigma}$ les éléments de structure caractéristiques des groupes de Galois arithmétiques -- et faire le lien avec les constructions typiques de cet ordre en théorie des corps de classes.

d) Examiner dans quelle mesure dans ce contexte on peut faire intervenir les conditions de compatibilité d'apparence draconiennes sur $\alpha,\beta$ du §47 liées à la compatibilité avec
\[
\diagram{\begin{tikzcd}[ampersand replacement=\&,column sep=large]
\varphi:\widehat\Pi'_{0,3} \arrow[r,"\sim"] \& \mathcal G_{0,3}^{+\wedge}
\end{tikzcd}}
\]
-- je \unsure{prévois.} d'ailleurs que le cadre des calculs précédents est trop étroit pour arriver à y faire intervenir ces conditions supplémentaires.

e) Ce qui m'intéresse le plus pour le moment, c'est une systématisation des calculs précédents dans un cadre beaucoup plus vaste que celui du seul groupe $SG$ et de l'hom. \unsure{standard}

\mspage{487}

\scanpage{489}
\mspage{547}

\[
\diagram{\begin{tikzcd}[ampersand replacement=\&,column sep=large]
\mathcal G_{0,3}^{+\wedge} \arrow[r] \& \deleted{\unsure{S}}G(\widehat{\mathbb Z})/\pm1
\end{tikzcd}}
\]
-- j'ai le sentiment d'avoir mis le doigt sur un procédé systématique pour construire des représentations ``motiviques'' de $\Gamma_{\mathbb Q}$, via des représentations de la partie ``géométrique'' $\mathcal G_{0,3}^{+\wedge}$ de
\[
\mathcal M_{0,3}^{+}\simeq\Pi_1(M_{0,3\,\mathbb Q},\unsure{\mathfrak S}_3)
\simeq\Pi_1(M'_{1,1\,\mathbb Q})
\]
-- ou mieux, \inserted{via} des représentations du groupe $\mathcal G_{0,3}^{+}\simeq\mathrm{Sl}(2,\mathbb Z)'$ discret (ou les \unsure{``}variantes telles $\mathrm{Sl}(2,\mathbb Z)$, $\Pi_{0,3}$ \ldots) à valeurs dans les points entiers de schémas en groupes (\inserted{\unsure{génériquement}} réductifs ??) sur $\mathbb Z$.

f) Il y a pourtant une autre direction d'investigation qui m'attire beaucoup, pas nécessairement liée à l'introduction des schémas en groupes, et qui pourrait se poursuivre \inserted{\unsure{en outre}} dans le cadre des seuls ``gros'' groupes fondamentaux profinis --

\mspage{488}

\scanpage{490}

savoir l'étude de certains \inserted{$\mathbb Q$-}revêtements particuliers de $\unsure{U}_{0,3\,\mathbb Q}=\mathbb P_{\unsure{0,3}}-\{0,1,\infty\}$, et des renseignements qu'on doit en tirer sur des propriétés supplémentaires des actions ext. ``arithmétiques'' de $\Pi_{0,3}^{\wedge}$ (du type des résultats obtenus dans §47), et \deleted{\unsure{sur}} des calculs de scindages de l'extension $\mathcal M_{0,3}$ de $\Gamma_{0,3}$ par $\mathcal G_{0,3}^{+\wedge}$ en d'autres points que les points ``critiques'' $0,1,\infty,\frac12,2,-1,-\bar j,\unsure{-j}$.

(Sans compter la question \unsure{d'asseoir|d'assurer} les fondements, via la méthode suggérée par Deligne -- et d'étudier l'équation des lacets dans le compactifié \unsure{protilp.|pro. tilp.} $\unsure{L_A|L_\Delta}$ -- et de continuer à investiguer dans la direction des conjectures \unsure{fondamentales}, notamment sur l'\unsure{arithmétisation} des points \uline{complexes} -- etc etc).

\mspage{489}

\scanpage{491}
\mspage{548}

Revenons pour l'instant à la question d'identifier les \inserted{quatre} hom. obtenus (33) à (36), en termes de choses classiques. On peut considérer l'hom. (33) \inserted{$\Theta_{B'_0}^{\pi}$} comme de nature essentiellement technique, ou du moins comme étant subordonné à l'hom. $\Theta_G^\pi$ de (34). D'ailleurs, (35) et (36) s'en déduisent, en \inserted{les} composant avec les deux hom.
\[
\diagram{\begin{tikzcd}[ampersand replacement=\&,column sep=large]
\Gamma_{0,3}^{\sim} \arrow[r] \& \mathcal M_{0,3}^{\sim}
\end{tikzcd}}
\]
\inserted{(moralement du moins, pour le moment !)} associés aux \deleted{\unsure{les}} points $\frac12$ et $-\bar j$ de $\unsure{U}_{0,3\,\mathbb Q}(\mathbb C)$, et qui correspondent, au niveau des schémas en groupes, aux deux scindages $\psi_\rho,\psi_\sigma$ de $\overset{\inserted{(3)}}{\mathcal M_{\rho,\sigma}}$ sur $\overset{\inserted{(31)}}{\mathcal T_{\rho,\sigma}}$, donnés \deleted{par} respectivement par
\diagram{\begin{equation*}
\tag{48}
\left\{\begin{gathered}
\begin{tikzcd}[ampersand replacement=\&,row sep=small,column sep=large]
\psi_\rho:(\underline a,\underline b) \arrow[r,mapsto] \& (\underline a,\underline b,\underline a) \\
\psi_\sigma:(\underline a,\underline b) \arrow[r,mapsto] \& (\underline a,\underline b,\underline b)
\end{tikzcd}
\end{gathered}\right.
\end{equation*}}
\margin{\diagram{\begin{tikzcd}[ampersand replacement=\&,row sep=small,column sep=large]
\deleted{\unsure{\mathcal T_{\rho,\sigma}|\psi_\rho}}\,\mathcal T'_{\rho,\sigma}\arrow[r] \& \mathcal M_{\rho,\sigma} \\
\mathcal T_{\rho,\sigma}\arrow[r] \& \mathcal M_{\rho,\sigma},
\end{tikzcd}}}
avec ce grain de sel que le premier n'est défini que sur le sous-groupe $\mathcal T'_{\rho,\sigma}$ défini par
\[
\tag{49}
\left\{\begin{aligned}
\mathcal T'_{\rho,\sigma}&\subset\mathcal T_{\rho,\sigma}\\
\Vert\quad&\left\{\left(\underset{\inserted{\substack{\rotatebox[origin=c]{-90}{$\unsure{\in}$}\\T_\rho}}}{\underline a},\underset{\inserted{N'_\sigma}}{b}\right)\,\middle|\,\deleted{\text{dans }\mathcal T_{\rho,\sigma}}\ \det\underline a=\det\underline b\right\}\\
&\Vert\quad T_\rho\mathbin{\underset{\mathbb{G}_m}{\times}}N'_\sigma.
\end{aligned}\right.
\]

\mspage{490}
